\section{Offenlegung, Impressum, Nutzungsbedingungen und Hinweise}

\begin{WidePar}

\begin{multicols}{2}
\minisec{Medieninhaber}

{\noindent\sffamily Arbeitsgemeinschaft Arbeits- und Ausbildungsstandards für
den Sanitätsdienst -- }{\ARGEAASS}
\\
\noindent{}Engerthstraße 146\,/\,8\,/\,13 \\
A-1200 Wien -- AUSTRIA

\begin{compactdesc}
  \item[E-Mail:] office@aass.at
  \item[Homepage:] \url{http://www.aass.at}
  \item[ZVR:] 846678982
\end{compactdesc}


\minisec{Weitere Informationen}

\begin{compactdesc}
  \item[Verlags- und Herstellungsort:] Wien
  \item[Publikationsart:] Elektronisch. Die etwaige Herstellung von
  Druckwerken erfolgt durch den autorisierten Nutzer.
  \item[ISBN:] ---
  \item[ISSN:] ---
  \item[Status:] \DocVersionTyp{} (\DocVersionInfo{})
  \item[Versionsgeschichte:]~
  
  {\scriptsize %
  \begin{description}
    \item[0.2.0] 2009-07 (VV\footnote{Vorabversion}, 349 Änderungen\footnote{Seit
    Übernahme in das Versionsverwaltungssystem.})
    \item[0.4.0] 2009-09 (VV, 284 Änderungen)
    \item[0.6.0] 2009-10 (VV, 322 Änderungen)
    \item[0.8.0] 2010-01 (VV, 653 Änderungen)
    \item[0.10.0] 2010-06 (VV, 935 Änderungen)
    \item[0.12.0] 2010-08 (VV, 489 Änderungen)
    \item[0.14.0] 2011-01 (VV, 706 Änderungen)
    \item[0.16.0] 2011-03 (VV, 141 Änderungen)
    \item[0.18.0] 2011-08 (VV, 449 Änderungen)
    \item[0.20.0] 2011-?? (VV, ??? Änderungen\footnote{Lt.
    vorläufiger Statistik, Änderungen nach Redaktionsschluss möglich.})
  \end{description}
  }
\end{compactdesc}

\noindent{}Die {\ARGEAASS} ist ein gemeinnütziger Verein mit Sitz in Wien,
welcher die Entwicklung, Wartung, Qualitätssicherung, Anpassung, Verbreitung und Lehre von
wissenschaftlich und sachlich fundierten Arbeits- und Ausbildungsstandards für
den Sanitätsdienst (AASS), sowie die wissenschaftliche Forschungstätigkeit in
diesem Gebiet bezweckt.

\minisec{Mitarbeit und Meldung von Fehlern}

\noindent{}Wünsche, Anregungen, Beschwerden und Fehlermeldungen nehmen wir
gerne unter 
% \url{http://www.aass.at/bugs}
\href{mailto:anregungen@aass.at}{\nolinkurl{anregungen@aass.at}} entgegen.

\minisec{Nutzungsbestimmungen}

\noindent{}Dieses Werk ist nur für den internen Gebrauch gem. §\,42
Abs.\,6 UrheberG bestimmt. Die Weitergabe ist nicht
gestattet. 

Sollten Sie Interesse an den Arbeits- und
Ausbildungsstandards für den Sanitätsdienst haben, setzen
Sie sich bitte mit uns in Verbindung.

\minisec{Hinweise}

\noindent{}Wie jede Wissenschaft ist die Medizin ständigen Entwicklungen
unterworfen. Soweit fachliche Angaben gemacht werden, darf der Leser
zwar darauf vertrauen, dass die Autoren große Sorgfalt
darauf verwandt haben, dass diese Angaben dem Wissensstand bei
Fertigstellung des Werkes entspricht. Für Angaben, speziell
über Dosierungsanweisungen und Applikationsformen, kann jedoch
keine Gewähr übernommen werden. 

Jeder Leser ist angehalten, durch
sorgfältige Prüfung der Literatur, der Lehrmeinung so\-wie der
Beipackzettel der verwendeten Präparate und gegebenenfalls nach
Konsultation eines Spezialisten festzustellen, ob die gegebene
Aussage oder Empfehlung für Dosierungen oder die Beachtung von
Kontraindikationen gegenüber der Angabe in diesem
Werk abweicht. Jede Dosierung oder Applikation erfolgt
auf eigene Gefahr des Benutzers. 

Die {\ARGEAASS} bittet jeden Leser, ihm aufgefallene Ungenauigkeiten und
Fehler, sowie Anregungen und konstruktive Kritik, über die Email-Adresse
\href{mailto:anregungen@aass.at}{\nolinkurl{anregungen@aass.at}} mitzuteilen.

Geschützte Warennamen und Warenzeichen werden nicht durchgehend besonders
kenntlich gemacht. Aus dem Fehlen eines solchen Hinweises kann also
nicht geschlossen werden, dass es sich um einen freien Warennamen
handelt.

\minisec{Die Technik}

\begin{compactdesc}
  \item[Satz und Druckvorstufe] \LaTeX{} mit {\sffamily KOMA-Script}
  
  Distribution: \TeX live 2009, u.\,a. mit den Paketen: AASS
  Style Collection (asc),
  \Code{microtype},
  \Code{babel},
  \Code{babelbib} (mehrsprachiges Literaturverzeichnis),
  \Code{shorttoc} (Inhaltsübersicht),
  \Code{titeltoc},
  \Code{makeidx} (Index),
  \Code{prettyref},
  \Code{units},
  \Code{nicefrac},
  \Code{soul},
  \Code{soulutf8},
  \Code{booktabs},
  \Code{marvosym} (Symbole),
  \Code{txfonts} (Symbole),
  \Code{skul} (Symbole),
  \Code{ifsym} (Symbole),
  \Code{dingbat} (Symbole),
  \Code{paralist},
  \Code{geometry} (Seitenlayout),
  \Code{fancybox},
  \Code{lmodern} (Latin Modern Fonts),
  \Code{titlesec} (Sectioning titles),
%   \Code{float},
  \Code{hyperref} (Hypertext links),
  \Code{url} (URL's),
  \Code{quotchap},
  \Code{minitoc} (Inhaltsverzeichnisse für Kapitel),
  \Code{chappg} (Seitennummerierung pro Kapitel),
  \Code{caption} (Customizing captions),
  \Code{caption3},
  \Code{mhchem} (chemische Formeln),
  \Code{optional} (optionaler Text),
  \Code{ulem}
  \item[Entwicklungsumgebung] Eclipse Ganymed mit Texlipse- und JGit/EGit-Plugin,
  Texmaker
  \item[Versionsverwaltung] Git
  \item[Literaturverzeichnis] JabRef 2.5.0, BibTeX, Babelbib
  \item[Textverarbeitung und Rohentwurf] OpenOffice
  \item[Betriebssysteme] \E{Rohentwurf}: Linux (Debian, Gentoo, Ubuntu), MacOS X,
  MS Windows XP, Vista
  
  \E{Endproduktion}: Ubuntu Linux 10.10 \textit{"`Maverick
  Meerkat"'}, \url{http://www.ubuntu.com}
\end{compactdesc}





% {\tiny
% scrkbase  v3.04a KOMA-Script package,
% scrbase   v3.04a KOMA-Script package,
% keyval   key=value parser (DPC),
% scrlfile KOMA-Script package (loading files),
% tocbasic  KOMA-Script package (handling toc-files),
% typearea a KOMA-Script package (type area),
% scrhack KOMA-Script package (hacking other packages),
% inputenc  Input encoding file,
% fontenc Standard LaTeX package,
% blindtext blindtext-Package,
% xspace Space after command names (DPC,MH),
% microtype  Micro-typography with pdfTeX (RS),
% fixltx2e 2006/09/13 v1.1m fixes to LaTeX,
% ellipsis 2004/9/28 v1.6 ellipsis: fixes spacing, 
% graphicx 1999/02/16 v1.0f Enhanced LaTeX Graphics (DPC,SPQR),
% graphics 2009/02/05 v1.0o Standard LaTeX Graphics (DPC,SPQR),
% trig 1999/03/16 v1.09 sin cos tan (DPC),
% babel 2008/07/06 v3.8l The Babel package,
% babelbib 2009/10/10 v1.29 babelbib: multilingual bibliographies (HH),
% ifthen 2001/05/26 v1.1c Standard LaTeX ifthen package (DPC),
% shorttoc 2002/08/20 v1.3 Short table of contents (JPFD),
% titletoc 2007/08/12 v1.6 TOC entries,
% makeidx 2000/03/29 v1.0m Standard LaTeX package,
% prettyref 1998/07/09 v3.0,
% aass-common,
% xcolor 2007/01/21 v2.11 LaTeX color extensions (UK),
% amsmath 2000/07/18 v2.13 AMS math features,
% amstext 2000/06/29 v2.01,
% amsbsy 1999/11/29 v1.2d,
% amsopn 1999/12/14 v2.01 operator names,
% amssymb 2009/06/22 v3.00,
% amsfonts 2009/06/22 v3.00 Basic AMSFonts support,
% units 1998/08/04 v0.9b Typesetting units,
% nicefrac 1998/08/04 v0.9b Nice fractions,
% soul 2003/11/17 v2.4 letterspacing/underlining (mf),
% soulutf8 2007/09/09 v1.0 Adding support for UTF-8 to soul (HO),
% infwarerr 2007/09/09 v1.2 Providing info/warning/message (HO),
% etexcmds 2007/12/12 v1.2 Prefix for e-TeX command names (HO),
% chngpage 2009/10/20 v1.2b change page layout,
% calc 2007/08/22 v4.3 Infix arithmetic (KKT,FJ),
% booktabs 2005/04/14 v1.61803 publication quality tables,
% longtable 2004/02/01 v4.11 Multi-page Table package (DPC),
% supertabular 2004/02/20 v4.1e the supertabular environment,
% tabularx 1999/01/07 v2.07 `tabularx' package (DPC),
% array 2008/09/09 v2.4c Tabular extension package (FMi),
% tabulary 2008/12/01 v0.9 tabulary package (DPC),
% multicol 2008/12/05 v1.6h multicolumn formatting (FMi),
% hhline 1994/05/23 v2.03 Table rule package (DPC),
% textcomp 2005/09/27 v1.99g Standard LaTeX package,
% marvosym 2006/05/11 v2.1 Martin Vogel's Symbols font definitions,
% txfonts 2008/01/22 v3.2.1,
% skull 2002/01/23 v0.1 (c) Henrik Christian Grove <grove@math.ku.dk>,
% savesymbol 2003/06/01 v1.1 Saves and restores symbols,
% ifsym 2000/04/18 I.Kloeckl,
% dingbat 2001/04/27 v1.00 Hands and other dingbats,
% paralist 2002/03/18 v2.3b Extended list environments (BS),
% marginnote 2009/05/06 v1.1d non floating margin notes for LaTeX,
% geometry 2010/03/01 v5.2 Page Geometry,
% ifpdf 2009/04/10 v2.0 Provides the ifpdf switch (HO),
% ifvtex 2008/11/04 v1.4 Switches for detecting VTeX and its modes (HO),
% layout 2000/09/25 v1.2c Show layout parameters,
% fancybox 2000/09/19 1.3,
% lmodern 2009/10/30 v1.6 Latin Modern Fonts,
% ocr 2006/09/18 LaTeX support for the various OCR fonts. Created by Pal,
% titlesec 2007/08/12 v2.8 Sectioning titles,
% wrapfig 2003/01/31  v 3.6,
% float 2001/11/08 v1.3d Float enhancements (AL),
% hyperref 2009/10/09 v6.79a Hypertext links for LaTeX,
% ifxetex 2009/01/23 v0.5 Provides ifxetex conditional,
% hycolor 2009/10/02 v1.5 Code for color options of hyperref/bookmark,
% xcolor-patch 2009/10/02 xcolor patch,
% kvoptions 2009/08/13 v3.4 Keyval support for LaTeX options (HO),
% kvsetkeys 2009/07/30 v1.5 Key value parser with default handler support,
% url 2006/04/12  ver 3.3  Verb mode for urls, etc.,
% bitset 2007/09/28 v1.0 Data type bit set (HO),
% intcalc 2007/09/27 v1.1 Expandable integer calculations (HO),
% bigintcalc 2007/11/11 v1.1 Expandable big integer calculations (HO),
% pdftexcmds 2009/09/23 v0.6 LuaTeX support for pdfTeX utility functions,
% ifluatex 2009/04/17 v1.2 Provides the ifluatex switch (HO),
% ltxcmds 2009/08/05 v1.0 Some LaTeX kernel commands for general use,
% atbegshi 2008/07/31 v1.9 At begin shipout hook (HO),
% quotchap 1998/02/09 v0.9f Decorative Chapter Headings with Quotes,
% minitoc 2008/07/16 v60 Package minitoc (JPFD),
% mtcmess 2006/03/14,
% chappg 2006/05/09 v2.1b page numbering by chapter number pages by chapter,
% caption 2009/10/09 v3.1k Customizing captions (AR),
% caption3 2009/10/09 v3.1k caption3 kernel (AR),
% ltcaption 2008/03/28 v1.2 longtable captions (AR),
% cite 2009/08/29  v 5.2,
% mhchem 2007/05/19 v3.07 for typesetting chemical formulae,
% twoopt 2008/08/11 v1.5 Definitions with two optional arguments (HO),
% excludeonly 2003/03/14 v1.0 eponymous command opposite to includeonly,
% optional 2005/01/26 ver 2.2b;  Optional inclusion/omission,
% glossaries 2010/02/06 v2.05 (NLCT),
% xkeyval 2008/08/13 v2.6a package option processing (HA),
% mfirstuc 2009/11/03 v1.04 (NLCT),
% xfor 2009/02/05 v1.05 (NLCT),
% translator 2007/03/11 ver 1.00,
% glossary-hypernav 2007/07/04 v1.01 (NLCT),
% glossary-list 2009/05/30 v2.01 (NLCT),
% glossary-tree 2009/01/14 v1.01 (NLCT),
% ulem 2000/05/26,
% nameref 2007/05/29 v2.31 Cross-referencing by name of section,
% refcount 2008/08/11 v3.1 Data extraction from references (HO),
% }
\end{multicols}
\end{WidePar}

% \noindent{}{\scriptsize \textbf{Vorabversionen:} 0.2.0 2009-07 / 0.4.0 2009-09
% / 0.6.0 2009-10 / 0.8.0 2010-01 / 0.10.0 2010-06}

% \paragraph*{Mission Statement} \textperiodcentered Das Ausbildungszentrum (ABZ) ist
% eine Dienst\-leistungs\-ein\-richtung des ASB Floridsdorf-Donau\-stadt
% und orientiert sein Angebot an den Be\-d\"urf\-nissen unserer Partner
% und an den strategischen Zielen des \T{ASB}. Wir sind sowohl f\"ur die 
% medizinische und sanitätsdienstliche Ausbildung unserer Mitarbeiter
% als auch f\"ur die Durchf\"uhrung von Schulungen und Lehrgängen zur
% Schulung der Bev\"olkerung verantwortlich. In Zusammenarbeit mit der
% \"Arztlichen Leitung und den \"ubergeordneten Gremien  entwickeln wir
% Arbeitsstandards und Lehrmeinungen. ü 